\SourceLeaf{379}%
hung mit mehr Vorsicht zu Werke. Sie dienten, durch Pulversignale auf dem Mont Tabor den Mont Cenis mit dem Mont Colombier, und durch Pulversignale auf dem Berge Pierre sur autre den Mont Colombier mit dem französischen Dreieckspunkte Puy d’Usson zu verbinden; zugleich wurden noch auf dem Mont Colombier selbst Signale gegeben, die zur Verknüpfung dieses Platzes mit der Sternwarte von Genf dienten. Die Zeitbestimmung auf dem Mont Cenis und dem Mont-Colombier war auf Beobachtungen an Mittagsfernröhren von \textsc{Lenoir}\footnote{Nach einigen Umständen zu schliessen, scheint 1821 und 1822 dasselbe Mittagsfernrohr auf dem Mont Cenis gebraucht zu sein, obgleich hier ein anderer Verfertiger genannt ist.} und \textsc{Grindel}, die für den französischen Standpunkt auf absolute mit einem Repetitionskreise gemessene Sternhöhen gegründet (vgl. \textit{Conn. des tems} 1829 und unsere Anz. 1828 Jan. 10). Auf dem Mont Cenis war man genöthigt, sich drei Stunden Weges von dem Hospiz jedesmal zu entfernen, um die Signale sehen zu können. Endlich finden wir hier noch die Bestimmung des Längenunterschiedes zwischen den Sternwarten von Turin und Mailand durch Pulversignale auf dem S. Bernardo di Fenera zu drei verschiedenen Zeiten \(1823\mathpunct{.}\mathpunct{.}\mathpunct{.}1824\), wobei alle Umstände so günstig waren, wie sie nur bei Operationen dieser Art sein können.

Der dritte Abschnitt enthält die Breitenbestimmungen der Sternwarten auf dem Mont Cenis, dem Mont Colombier und in Turin, die beiden ersteren mit Repetitionskreisen von \textsc{Troughton} und \textsc{Reichenbach}, die letzte mit dem \textsc{Reichenbach}schen Meridiankreise. Die letztern Beobachtungen zeigen nicht ganz den Grad von Übereinstimmung, an welchen man sonst bei diesen Instrumenten gewöhnt ist. Die Verf. haben dies selbst bemerklich gemacht. und lassen es auf sich beruhen, ob solche Anomalien Realität haben, oder von irgend einem Fehler in der Behandlung des Instruments abhängen, der sich in Zukunft aufklären lassen werde. Ref. bescheidet sich, dass bei der Mannigfaltigkeit der Aufmerksamkeiten, welche dieses Instrument erfordert, niemand, ohne an Ort und Stelle zu sein, auch nur eine plausible Vermuthung darüber aufstellen könne, findet aber in der Art, wie die Verf. sich über jene Anomalien geäussert haben, eine Aufforderung, aus seiner eigenen Erfahrung ein Beispiel anzuführen, wie geringfügige Umstände zuweilen den Beobachtungen nachtheilig werden können. Während einer Reihe von Jahren war die schöne Harmonie in den Beobachtungen an einem dem Turiner ganz gleichen Meridiankreise nur ein einzigesmal eine Zeit%
